



\section{Number Theory}
\begin{state}{ }
    Let $p \in \mathbb{Z} with p > I.$ Then p is said to be a prime if the 
    only positivive divisors of $p$ are $1$ and as $p$.

    \textbf{Remark}: The positive integer 1 is neither prime nor composite.

\end{state}

\begin{state}{Lemma 1.5}
    Every integer greater than $1$ has a prime divisor.
\end{state}
\subsection{}
\begin{subprove}
    Assume there exists an integer $n>1$ that has no prime divisor.
    By well-ordering principle \footnote{For a finite set of positive integers, there is always lower bound.}, we nat assume that $n$ is the least such integer.
    Clearly, $n~|~n$, so $n$ is not prime. Then, $n$ is composite and consequently there exist.
\end{subprove}

$a,b\in \mathbb{Z}$ such that $n=ab$, where $1<a<n$ and $1<b<n$. Since $1<a<n$, we have that $a$ has a 
prime divisor, say $P$, so that $P~|~a$. But, $a|n$, so we have that $p|n$ by proposition 1.1 
(transitivity of divisibility). Contradictionan, thus, Every integer greater than $1$ has a prime divisor.





\begin{state}{Theorem 1.6: Euclid}
    There are inifinitely many primes.
\end{state}
\begin{subprove}
    Assume that there is only finitely many primes where ${p_1,p_2, \cdots}$.
    Consider the integer $N=p_1\cdot p_2 \cdots p_n +1$, by lemma 1.5, we 
    know that $N$ has a prime divisor, say $p$.

    Thus $p=p_i$ for some $i=1,2,\cdots,n$. Then, $p|(N-p_1\cdot p_2\cdots p_n)$ by 
    proposition 1.2. So $p|1$, a contradition.

    Thus there are infinitely many prime numbers.

\end{subprove}

\begin{state}{prop 1.7}
    Let $n$ be a composite number. Then, n has a prime divisor $p$ with $p\leq \sqrt{n}$
\end{state}
\begin{subprove}
    Since n is composite, there exists $a,b\in\mathbb{Z}$ such that $n=ab$ with $1<a<n$ and $1<b<n$, 
    and without loss of generality, $a\leq b$, we must have $a \leq \sqrt{n}$, otherwise, $a>\sqrt{n}$ 
    then $n=ab>\sqrt{n}\sqrt{n}=n$. Contradiction, by lemma 1.5, 
    we know that $a$ must have a prime divisor, say $p$. Since $p|a$ and $a|n$ we have that $p|n$ by prop 1.1.
    Furthermore, $p \leq a \leq \sqrt{n}$, so $p$ is the desired prime factor of $n$. 
\end{subprove}

If you do larger experiments with the sieve of eratosthenes one might 
guess that primes are 
more sparsely distributed as you progress to larger and larger integers.


\begin{state} {Proposition 1.8}
    For any positive integer $n$, there are at least $n$ consecutive composite positive integers.
\end{state}
\begin{subprove}
    Given $n\geq 1$, consider the sequence
    $$(n+1)! + 2, (n+1)! + 3,\cdots,(n+1)! + n+1$$
    Let $i$ be an integer such that $2\leq i \leq n+1$ 
    Since $i|(n+1)!$ and $i|i$, we have $i|((n+1)!+i)$ for $2\leq i \leq n+1$ by proposition 1.2.
    Thus, each integer in the sequence is composite.
\end{subprove}

\begin{state}{ }
    Any pair $\left(p,p+2\right)$ with both entries prim are called twin prime. (not far, but close to each other.)
\end{state}

\begin{state}{conjecture}
    There are infinitely many primes $p$ such that $p+2$ is also prime.
\end{state}

Recent breakthroughs: Zhang (2013): Let $p_1, p_2, \cdots$, denote the sequence of primes. 
Zhang proved that there are inifinitely many primes $p_n$ such that
$$p_{n+1}-p_{n} \leq 70~000~000$$

Improved by Maynard, Tao, ..., but not all the way to 2.


\begin{state}{Prime Number Theorem (1793) \footnote{Conjectured by Gauss and proved in 1896 by de la Vallee Poussin and Hadamard.}}
    Let $x\in \mathbb{R}$ with $x>0$, then $\Pi (x)=\left|{1<p\leq x}\right|$ where
    $$\lim_{x\to \infty}{\frac{\Pi (x)}{\frac{x}{\log{x}}}} = 1$$
\end{state}

Many unsolved problems in number theory deal with integers that are expressible in certain forms.
\begin{state}{Goldbach Conjecture (1742)}
    Every even integer greater than $2$ can be expressed as the 
    sum of two (not necessarily distinct) primes. 
\end{state}

Hardy and Littlewood predicted a formula for the number of 
representations of an even integer has in terms of a sum of two prime.

\begin{state}{Definition}
    any prime number expressible in the form $2^p-1$ with $p$ prime is 
    siad to be a Mersenne prime.
\end{state}

\begin{state}{Conjecture}
    There are inifinitely many Mersenne primes.
\end{state}

\begin{state}{Conjecture}
    There are infinitely many primes of the form $n^2+1$.
\end{state}
